\section*{الفصل \ref{ch:g9:evidence-of-evolution} --- أدلة التطور}

\begin{solution}{exo:g9:evidence-of-evolution:1}
أثر محفوظ من حياة ماضية مختوم في صخر رسوبي. والطبقات
تتراكم بترتيب الترسب، أقدمها أسفلها --- فالصعود في
الكومة تقدم في الزمن.
\end{solution}

\begin{solution}{exo:g9:evidence-of-evolution:2}
\omterm{def:g1:animals-around-us:animal}{الحيوانات} تتغير طبقة بعد طبقة، وأكثر \omterm{def:g5:biodiversity:species}{أنواع} العمق
منقرضة؛ والمجموعات الكبرى تظهر بترتيب تعشش الشجرة ---
\omterm{prop:g3:sorting-living-things:groups}{الأسماك} في العمق، وأقرباء ذوات \omterm{def:g1:my-body:parts}{الأطراف} الأربعة فوقها،
\omterm{prop:g3:sorting-living-things:groups}{والثدييات} \omterm{prop:g3:sorting-living-things:groups}{والطيور} في الأعلى؛ والوسائط تقف عند نقاط
التفرع.
\end{solution}

\begin{solution}{exo:g9:evidence-of-evolution:3}
\omterm{def:g1:animals-around-us:covering}{ريش} وجناحان --- عدة الطائر --- على هيكل بأسنان وذنب
عظمي طويل --- عدة الزاحف: المجموعتان في \omterm{def:g1:animals-around-us:animal}{حيوان} واحد، عند
الفرق الذي ترسمه الشجرة.
\end{solution}

\begin{solution}{exo:g9:evidence-of-evolution:4}
\omterm{def:g3:bones-and-muscles:skeleton}{عظام} ورك \omterm{def:g1:my-body:parts}{ورجل} الحوت المخفية: انحدار ماش. وجناحا الخنفساء
غير الطائرة الملحومان: أسلاف طائرة. وعصعص الإنسان (أو
الزائدة الضامرة): خط ذو ذنب، مختلف التغذية. وكل واحد
منها جزء أبقي يشهد ببناء سابق.
\end{solution}

\begin{solution}{exo:g9:evidence-of-evolution:5}
عدّ فروق أحرف مورّثة مشتركة بين نوعين: فكلما قلّت قربت
\omterm{prop:g6:classification-kinship:kinship}{القرابة}. والمسافات، إذا جمعت، ركّبت شجرة --- وهي تطابق
التي رسمتها \omterm{def:g3:bones-and-muscles:skeleton}{العظام} \omterm{def:g9:evidence-of-evolution:fossils}{والأحافير}.
\end{solution}

\begin{solution}{exo:g9:evidence-of-evolution:6}
الصورة الداكنة من عثة الفلفل تتبع السخام وصفاء الهواء؛
وآفات منيعة على المبيدات والمضادات الحيوية تنشأ في
سنوات؛ وقياسات مناقير عصافير الجزر تنزاح مع الجفاف
والوفرة.
\end{solution}

\begin{solution}{exo:g9:evidence-of-evolution:7}
هندسة جديدة كانت تستطيع أن تعطي كل طرف آليته المثلى ---
دعامات للأجنحة، وأجنحة مائية للزعانف --- وما كانت لتكرر
عدد \omterm{def:g3:bones-and-muscles:skeleton}{عظام} اعتباطيًا. أما الترتيب الواحد، ممطوطًا بعناء
إلى كل حرفة، فهو ما ينتجه البناء على الموروث ولا ينتجه
تصميم حر.
\end{solution}

\begin{solution}{exo:g9:evidence-of-evolution:8}
الطائفة: \omterm{def:g2:animals-grow-up:born}{الحليب}، \omterm{def:g7:what-is-respiration:respiration}{وتنفس} الهواء، والشعر الجنيني ---
موروثات \omterm{prop:g3:sorting-living-things:groups}{الثدييات}
(\cref{ex:g6:classification-kinship:whale}). والأعضاء
الأثرية: أوراك \omterm{def:g3:bones-and-muscles:skeleton}{وعظام} \omterm{def:g1:my-body:parts}{أرجل} مخفية. \omterm{def:g9:evidence-of-evolution:fossils}{والأحافير}: السلسلة
المرتبة زمنيًا من ثدي شاطئي رباعي \omterm{def:g1:my-body:parts}{الأرجل} إلى سابح بلا
\omterm{def:g1:my-body:parts}{أرجل}. والجزيئات: مسافات مورّثاته تضعه بين \omterm{prop:g3:sorting-living-things:groups}{الثدييات}،
وأقرب أقربائه ذوات الحوافر. أربعة شهود، وملف واحد.
\end{solution}

\begin{solution}{exo:g9:evidence-of-evolution:9}
لأن الشهود لا يمكن أن يتواطؤوا: فالصخور رسّبت، والأجسام
بنيت، والجزيئات هجّئت بعمليات لا سبيل لأي منها إلى شهادة
الأخرى. والشهود المستقلون إذا اتفقوا ضاعفوا المصداقية
--- وشجرة واحدة من ثلاثة مصادر هي التقارب الذي يعتز به
المنهج.
\end{solution}

\begin{solution}{exo:g9:evidence-of-evolution:10}
لأن الدعوى تحرّمه: \omterm{prop:g3:sorting-living-things:groups}{فالثدييات} يجب ألا تظهر تحت طبقات
أسلافها. وكل حفر فرصة جديدة للتكذيب --- والغياب، مطردًا
في كل تنقيب أجري قط، هو الدعوى تجتاز اختبارها يوميًا.
\end{solution}

\begin{solution}{exo:g9:evidence-of-evolution:11}
\omterm{def:g5:human-reproduction:pregnancy}{مضغ} \omterm{def:g4:classifying-animals:vertebrate}{الفقاريات} تتفق تشابهًا لافتًا خلال أسابيع وضع المخطط
الباكرة وتفترق متأخرة. وآلة فصل النمو تقول لماذا:
فالتعليمات الباكرة تضع مخطط الجسم المشترك --- وهي أعمق
فصول النص وأقدمها --- بينما تكتب القراءة المتأخرة
خصوصيات \omterm{def:g5:biodiversity:species}{النوع}. نص باكر مشترك، وانحدار عميق مشترك.
\end{solution}

\begin{solution}{exo:g9:evidence-of-evolution:12}
الدعوى: \omterm{def:g5:biodiversity:species}{الأنواع} تتغير وتشترك في الانحدار. والتنبؤات:
\omterm{def:g1:animals-around-us:animal}{حيوانات} مرتبة زمنيًا في تعشش الشجرة؛ ووسائط عند الفروق؛
وأعضاء أثرية من بناءات سابقة؛ ومسافات جزيئية تتدرج مع
\omterm{prop:g6:classification-kinship:kinship}{القرابة}. والشهود: الصخور تقدم الترتيب والوسائط؛ والأجسام
تقدم المخطط والأعضاء الأثرية؛ والجزيئات تقدم الشجرة
المطابقة؛ والمراقبات الحية تقدم التغير متحركًا. والحكم:
الانحدار مع التغير قائم على شهادة مستقلة متقاربة ---
ويسلّم الفصل التالي سؤاله: \emph{ما الذي يقود التغير؟}
\end{solution}

\begin{solution}{pb:g9:evidence-of-evolution:1}
\textbf{1.} الترتيب: ثلاثيات الفصوص أسفل، ثم أول سمكة،
وآثار أقدام أقرباء \omterm{prop:g3:sorting-living-things:groups}{البرمائيات}، \omterm{prop:g4:classifying-animals:classes}{والزواحف}،
و\emph{الأركيوبتيريكس}، وفك ثدي مبكر. والتعليق: ``مقروءة
صعودًا، تغير الصخور شخوصها بالترتيب الذي تتنبأ به \omterm{def:g6:classification-kinship:tree}{شجرة
العائلة} تمامًا --- والوسائط واقفة عند التفرعات.''

\textbf{2.} التعليق: ``مخطط عظمي واحد --- ممطوط ليطير،
ويجدّف، ويعدو، ويكتب: \omterm{def:g9:heredity-and-traits:heredity}{وراثة} أعيدت وظيفتها، ولم يعد
تصميمها قط.'' والخزانة: أوراك الحوت، وجناحا الخنفساء
الملحومان، وعصعص الإنسان.

\textbf{3.} الشجرة التي تقتضيها أعداد الأحرف ---
الشمبانزي أقربها، \omterm{ex:g6:producing-our-food:bread}{والخميرة} أبعدها --- يجب أن تطابق شجرة
التشريحيين وعلماء \omterm{def:g9:evidence-of-evolution:fossils}{الأحافير} في الغرفة المجاورة: التفرع
نفسه، مستنبطًا استقلالًا.

\textbf{4.} العثات: ألوان \omterm{def:g8:reproduction-and-environment:population}{جمهرة} تتبع بيئتها في غضون
عقود. والعصافير: توزيعات المناقير تتحرك مع مؤونة
\omterm{def:g2:from-seed-to-plant:seed}{البذور}، مقيسة بالفرجار --- تغير تغذيه \omterm{def:g9:heredity-and-traits:heredity}{الوراثة}، مشاهدًا
حيًا.

\textbf{5.} لأن الانحدار يتفرع: فنحن وقرود اليوم أقرباء
على فرعين منفصلين من فرق مشترك --- فما من حي هو سلف
لأحد، ووجود ابن العم لا ينفي \omterm{prop:g6:classification-kinship:kinship}{القرابة} كما لا ينفي أبناء
عمومتك قرابتك.

\textbf{6.} \emph{الأركيوبتيريكس}، واقفًا بين هيكل
الزاحف وعدة الطائر؛ والحيتان الأحفورية ذوات \omterm{def:g1:my-body:parts}{الأرجل}،
واقفة بين الثدي الشاطئي والسابح. وكلاهما كان تنبؤًا قبل
أن يكون اكتشافًا.

\textbf{7.} في العلم تسمي ``النظرية'' تفسيرًا اختبر حتى
صار ينظّم حقلًا --- لا حدسًا. والذي تثبته الغرف هو
الحقيقة التاريخية للتغير والانحدار، بشهود مستقلين
متقاربين؛ أما النظرية بمعناها الدقيق فهي الآلية التي
تشرحها القاعة التالية.

\textbf{8.} مفارقة تاريخية حقيقية --- \omterm{def:g3:bones-and-muscles:skeleton}{عظام} ثدييات تحت
طبقات ثلاثيات الفصوص، موثقة --- كانت ستكسر الترتيب الذي
تقتضيه الدعوى. وألا يكون حفر واحد في قرنين قد أنتج
واحدة، فذلك هو الدعوى تنجو من أقسى اختباراتها، كل يوم،
وفي كل مكان يستمر فيه الحفر.

\textbf{9.} ``العلب داخل العلب التي لقيتها في الغرفة
المجاورة لم تكن قط تيسيرًا في الحفظ: \omterm{def:g9:heredity-and-traits:traits}{فالصفات} تتعشش لأن
الانحدار يتفرع --- \omterm{def:g4:classifying-animals:classification}{فالتصنيف} كان \omterm{def:g6:classification-kinship:tree}{شجرة العائلة}، مرسومة
قبل أن يعرف أحد ما الذي يرسمه.''

\textbf{10.} ``كل معروض في هذا الجناح، وكل زائر يقف
أمامه، يكتب بناءه بالأحرف الأربعة نفسها تقرؤها الشفرة
نفسها: فالفروع تختلف، ولغة الجذع لم تختلف قط.''

\textbf{11.} ``يبين هذا الجناح أن \omterm{def:g5:biodiversity:species}{الأنواع} تتغير، وأنها
تشترك في الانحدار، وأن السجل يحفظ الوسائط. أما ما لم
يقله بعد فهو ما الذي \emph{يقود} التغير --- وقاعة
الآلية هي التالية، وجوابها يجري على التنوع \omterm{def:g9:heredity-and-traits:heredity}{والوراثة}
وكوابح \omterm{def:g6:exploring-our-environment:components}{البيئة}.''

\textbf{12.} ``الحياة كلها ألبوم عائلي واحد، والصخور
حفظت صفحاته الأولى --- وبمهلة من الزمن يصير أبناء
العمومة مختلفين كالحيتان والخفافيش، ويبقون عائلة.''
\end{solution}
